\chapter{CUTLASS and CuTe GEMM Operations}\label{ch:cutlass}

\emph{Placement in the stack.} The distributed study of
Chapter~\ref{ch:distributed} fixed the per-GPU memory envelope
under each parallelisation strategy. Inside that envelope, the
wall-clock cost of a training step is largely determined by a
small set of GEMM shapes that every transformer layer re-issues at
every step. The present chapter characterises those shapes and
shows how each of them can be driven, in CUTLASS/CuTe, by a kernel
that reaches close to the peak Tensor-Core utilisation of the
device. The epilogue-fusion primitives built here are picked up
later: by the attention-output writeback of Chapter~\ref{ch:swfa}
and by the logit-projection GEMM of Chapter~\ref{ch:sglang}. What
this chapter's numbers establish, then, is not only a kernel-level
result but also an upper bound on what those higher layers can
subsequently deliver.

Matrix multiplication is, and will remain, the dominant consumer of
runtime in a transformer-based LLM. Every projection layer, every
MLP block, and both halves of the attention computation reduce in
the end to a GEMM---sometimes batched, sometimes grouped, but
always a GEMM. The chapter that follows describes how
high-performance GEMM kernels are built on top of NVIDIA's CUTLASS
library and CuTe layout algebra; these two components together
constitute the kernel layer of the proposed software stack.

\section{Motivation}\label{sec:cutlass_motivation}
Although PyTorch routes most of its matrix multiplications through
cuBLAS, the LLM workloads that reach production depend on several
GEMM variants that cuBLAS does not expose directly:
\begin{itemize}
    \item \textbf{Grouped and variable-shape GEMMs} appear in
          Mixture-of-Experts layers, where each expert has its own
          $(M_e, N_e, K)$ shape; the launch scheduler then has to
          pack all experts into a single kernel invocation.
    \item \textbf{Quantised GEMMs} in INT8, FP8, or INT4 are
          required on the inference side and demand mixed-precision
          accumulators together with a dequantisation epilogue.
    \item \textbf{Epilogue fusion}---bias add, GELU, residual add,
          or LayerNorm folded into the writeback of the main
          GEMM---lowers memory traffic by removing the second pass
          that would otherwise be needed over the output tensor.
    \item \textbf{Fused attention} co-designs the softmax and the
          value-weighted sum with the score GEMM itself; Flash
          Attention~\cite{dao2022flashattention} is the canonical
          example.
\end{itemize}
CUTLASS~\cite{cutlass} is a C++ template library that lets such
kernels be written at the granularity of Tensor Cores, with CuTe
serving as the layout and tensor algebra that sits underneath it.
Between them they occupy a useful middle ground: more flexible
than cuBLAS (which is opaque, hard to extend, and essentially
impossible to fuse into) and safer than raw CUDA (which is easy to
misuse and tedious to retarget across architectures).

\section{The CUTLASS Programming Model}\label{sec:cutlass_model}
CUTLASS expresses a GEMM of the form
$C \leftarrow \alpha\, A B + \beta\, C$ as a nested hierarchy of
tiles that align with the execution hierarchy of the GPU itself:

\begin{center}
\begin{tabular}{ll}
\textbf{Level} & \textbf{GPU abstraction} \\
\toprule
Threadblock tile & Cooperative Thread Array (CTA) / shared memory \\
Warp tile        & Warp (32 threads) \\
Instruction tile & Tensor Core MMA (e.g., \texttt{mma.m16n8k16}) \\
\bottomrule
\end{tabular}
\end{center}

\begin{figure}[H]
    \centering
    \begin{tikzpicture}[
        lvl1/.style={draw, thick, fill=blue!5,  minimum width=60mm,
                     minimum height=45mm, rounded corners=2pt},
        lvl2/.style={draw, thick, fill=blue!12, minimum width=22mm,
                     minimum height=15mm, rounded corners=1pt},
        lvl3/.style={draw, fill=orange!70, minimum width=5.5mm,
                     minimum height=4.5mm, inner sep=0pt, font=\tiny},
        lbl/.style={font=\footnotesize}
    ]
    % Outer: threadblock (CTA)
    \node[lvl1]  (cta) at (0, 0) {};
    \node[lbl, anchor=north west] at (cta.north west)
         {\bfseries Threadblock tile $B_m\!\times\!B_n$};

    % Inner: 4 warps (warp tile)
    \foreach \x/\y/\name in {-15/6/W0, 15/6/W1, -15/-12/W2, 15/-12/W3} {
        \node[lvl2] (\name) at (\x mm, \y mm) {};
        \node[lbl]  at (\name.north)
             [anchor=south, yshift=-12mm] {Warp tile};
    }

    % Inside one warp: MMA instruction fragments 4x3
    \foreach \c in {0,...,3} {
        \foreach \r in {0,...,2} {
            \node[lvl3] at
                 ($(W0.center) + (\c*6mm - 9mm, \r*5mm - 4.5mm)$) {};
        }
    }
    \node[lbl, font=\scriptsize, anchor=west, xshift=14mm] at (W0.east)
         {$16{\times}8{\times}16$ MMA};

    % Hardware mapping labels on the right
    \node[lbl, anchor=west] at (3.6, 2.0)  {$\leftarrow$ CTA / shared memory};
    \node[lbl, anchor=west] at (3.6, 0.3)  {$\leftarrow$ Warp (32 threads)};
    \node[lbl, anchor=west] at (3.6,-1.4)  {$\leftarrow$ Tensor-Core MMA};
    \end{tikzpicture}
    \caption{CUTLASS three-level tile hierarchy on Ampere/Hopper. A
             threadblock tile lives in shared memory, is partitioned
             into four warp tiles, and each warp tile is covered by
             a grid of instruction-level MMA fragments that issue
             directly to the Tensor Cores.}
    \label{fig:cutlass_tile_hierarchy}
\end{figure}

For input matrices
$A \in \mathbb{R}^{M \times K}$ and
$B \in \mathbb{R}^{K \times N}$, the kernel proceeds as follows:
\begin{enumerate}
    \item The output $C$ is partitioned into threadblock tiles of
          size $\mathtt{Bm} \times \mathtt{Bn}$.
    \item Each CTA iterates along the $K$ dimension in chunks of
          size $\mathtt{Bk}$, cooperatively copying the matching
          $A$ and $B$ tiles into shared memory.
    \item Within the CTA, warps read from shared memory into
          registers and issue MMA instructions, accumulating the
          result into an FP32 or FP16 register tile.
    \item An \emph{epilogue} writes the final tile back to global
          memory, optionally folding in a bias, an activation, and
          a residual add along the way.
\end{enumerate}

\subsection{Tile Sizes on A100}\label{sec:cutlass_tile_sizes}
The choice of tile size is always a compromise between how much
data is reused inside the tile and how much register or shared
memory the tile consumes. On A100 the sizes that work well in
practice are:
\begin{center}
\begin{tabular}{lll}
\textbf{GEMM shape class}    & \textbf{Threadblock} & \textbf{Warp} \\
\toprule
Large square / MLP           & $128 \times 128 \times 32$ & $64 \times 64 \times 32$ \\
Skinny (attention score)     & $128 \times 64  \times 32$ & $64 \times 32 \times 32$ \\
Very skinny (decode attn)    & $128 \times 16  \times 32$ & $64 \times 16 \times 32$ \\
Grouped (MoE)                & $128 \times 128 \times 32$ & $64 \times 64 \times 32$ \\
\bottomrule
\end{tabular}
\end{center}
The warp tile is sized so that each warp holds an integer number
of $16 \times 8 \times 16$ MMA fragments in its registers without
ever spilling.

\section{Software Pipelining}\label{sec:cutlass_pipelining}
HBM latency is hidden by multi-stage \emph{software pipelining},
which has the following ingredients:
\begin{itemize}
    \item Both Ampere and Hopper expose asynchronous copy
          instructions---\texttt{cp.async} on Ampere and the TMA
          copy engine on Hopper---that begin to fetch the next
          $K$-tile while the current tile is still being consumed
          by MMA instructions.
    \item CUTLASS keeps at least three shared-memory buffers in
          flight so that the copy engine, the MMA unit, and the
          epilogue never have to wait on one another.
    \item Swizzled shared-memory layouts guarantee that the MMA
          load patterns land on banks without collision.
    \item \texttt{cp.async} commit/wait groups are surfaced
          through CUTLASS's \texttt{MmaMultistage} class, which
          enqueues $N_{\text{stage}}$ fetches ahead of the MMA
          that will eventually consume them.
\end{itemize}

The effective HBM latency visible to each GEMM step is
\begin{equation}
    \tau_{\text{eff}} \;=\; \max\bigl(\tau_{\text{HBM}},
            N_{\text{stage}} \cdot \tau_{\text{MMA}}\bigr),
\end{equation}
so with $N_{\text{stage}} = 3$ and an A100 $\tau_{\text{MMA}}$ of
roughly $2$ cycles per $16 \times 8 \times 16$ MMA, the pipeline
hides up to $\sim\!6$ cycles of latency per MMA---more than enough
to swallow an L2 hit or a shared-memory conflict. An HBM miss of
around $200$ cycles is, however, too long for the innermost
pipeline stage to hide, which is precisely why the prefetch is
launched one stage earlier.

\section{CuTe Layouts}\label{sec:cute_layouts}
In CuTe~\cite{cutlass} every tensor is modelled as a pair
\[
    T = (\text{shape},\text{stride}),
\]
where the \emph{shape} is a (possibly nested) tuple that describes
the coordinate space and the \emph{stride} is a tuple that
describes the map from coordinates to a flat linear index. A plain
row-major $M \times N$ matrix, for instance, has
shape $(M, N)$ and stride $(N, 1)$.

\subsection{Composition}\label{sec:cute_composition}
The key operation in the algebra is \texttt{Layout} composition.
Given two layouts $L_A$ and $L_B$, the composite $L_A \circ L_B$
reinterprets the coordinates of $L_B$ through $L_A$. A typical
example is composing a row-major layout of shape $(8, 8)$ with a
swizzle layout: the result permutes the columns in a way that
avoids $32$-way bank conflicts in shared memory. That single
primitive is enough to describe, concisely:
\begin{itemize}
    \item the tiled partitions of $A$ and $B$, in any of their
          row-major, column-major, blocked, or swizzled forms;
    \item the warp-level and thread-level partitions of a
          shared-memory tile;
    \item the register-level fragment shapes that each variant of
          the MMA instruction requires as input.
\end{itemize}

\subsection{Hierarchical Tiling}\label{sec:cute_tiling}
A single transformer GEMM can be written as a short chain of
\texttt{local\_tile} calls:
\begin{lstlisting}[language=C++,basicstyle=\ttfamily\scriptsize]
// A: (M, K) with layout mA, tile into (BLK_M, BLK_K) per CTA
auto gA = local_tile(mA, make_shape(BLK_M, BLK_K),
                     make_coord(blockIdx.x, _));
// Warp-level partition within the CTA tile
auto tAsA = thrd_partition_A(warp_mma, gA);
// Copy to shared memory with an async-copy layout
auto tAgA = copy(copy_atom_cp_async, tAgA, tAsA);
\end{lstlisting}
One benefit of this style is that a single kernel template can be
retargeted to a wide range of data types, tile shapes, and hardware
generations simply by changing its layout parameters---the actual
MMA instruction is selected downstream by the \texttt{Sm80} or
\texttt{Sm90} traits class.

\section{Mapping LLM Layers onto GEMM}\label{sec:cutlass_llm}
Consider a transformer layer with hidden size $d_{\text{model}}$,
intermediate size $d_{\text{ff}}$, batch size $B$, and sequence
length $N$. Its constituent GEMMs look as follows:
\begin{itemize}
    \item \textbf{QKV projection.} Three GEMMs of shape
          $(BN) \times d_{\text{model}} \cdot d_{\text{model}}
          \times d_{\text{head}}$, which in practice are fused into
          a single GEMM of shape
          $(BN) \times d_{\text{model}} \cdot d_{\text{model}}
          \times 3d_{\text{head}}$. A bias-add epilogue eliminates
          one extra HBM pass over the QKV tensor.
    \item \textbf{Attention score.} A batched GEMM of
          $Q \cdot K^{\!\top}$ of shape
          $N \times d_{\text{head}} \cdot d_{\text{head}}
          \times N$ per head. Because the shape is skinny in
          $d_{\text{head}}$ (typically either $64$ or $128$),
          cuBLAS leaves performance on the table here.
    \item \textbf{Attention output.} A batched GEMM of the
          softmaxed scores with $V$, of shape
          $N \times N \cdot N \times d_{\text{head}}$. Flash
          Attention fuses both attention GEMMs and the softmax into
          a single kernel driven entirely by CuTe tiles.
    \item \textbf{MLP.} Two GEMMs, of shapes
          $(BN) \times d_{\text{model}} \cdot d_{\text{model}}
          \times d_{\text{ff}}$
          and
          $(BN) \times d_{\text{ff}} \cdot d_{\text{ff}}
          \times d_{\text{model}}$. A GELU epilogue is fused into
          the first of the two.
    \item \textbf{MoE.} A grouped GEMM spanning $E$ experts. Each
          expert has its own $(M_e, N_e, K)$ shape, drawn from the
          routing histogram at the previous step. The
          \texttt{GemmGrouped} kernel packs all groups into a
          single launch and thereby amortises the launch overhead
          across the $E$ experts.
\end{itemize}
CUTLASS ships templated kernels that cover all of the above,
including \texttt{GemmGrouped} for MoE and
\texttt{AttentionKernel} for Flash-Attention-style fusion.

\section{Epilogue Fusion}\label{sec:cutlass_epilogue}
A large share of the kernel launches inside an LLM are in fact
element-wise operations---bias add, GELU, residual add, RMSNorm.
Taken in isolation every one of these is memory-bound. CUTLASS's
epilogue mechanism lets us fold them into the writeback of the
preceding GEMM, so that the overall operation reads
\begin{equation}
    C \;\leftarrow\; \phi\!\left(\alpha\, A B +
                                 \beta\, C_{\text{prev}} +
                                 \text{bias}\right),
\end{equation}
with $\phi$ an activation function such as GELU or SiLU. The fused
form removes one full HBM pass over $C$ and is typically worth
somewhere in the range of $10$ to $25\%$ of end-to-end throughput
on an A100.

\subsection{Fused Bias+GELU
Epilogue}\label{sec:cutlass_bias_gelu}
A bias plus GELU epilogue is expressed in CUTLASS as follows:
\begin{lstlisting}[language=C++,basicstyle=\ttfamily\scriptsize]
using Epilogue = cutlass::epilogue::thread::LinearCombinationGELU<
    ElementOutput,        // FP16
    128 / cutlass::sizeof_bits<ElementOutput>::value,
    ElementAccumulator,   // FP32
    ElementCompute>;      // FP32
\end{lstlisting}
The important point is that the combination
$\beta \cdot C_{\text{prev}} + \text{bias}$ is evaluated in FP32
inside the epilogue's register tile, \emph{before} the GELU is
applied, so that numerical accuracy is preserved even when the
output $C$ is ultimately stored in FP16.

\subsection{Impact on HBM
Bandwidth}\label{sec:cutlass_hbm_impact}
The HBM traffic saved by epilogue fusion for the MLP GEMMs of a
GPT-2 Medium layer at $N = 2048$ and batch~$4$ is reported in
Table~\ref{tab:epilogue_bw}. Fusion eliminates two HBM
round-trips per MLP block, which corresponds to roughly $1$\,GB of
bandwidth per layer per step.

\begin{table}[h]
\centering
\caption{HBM traffic with and without epilogue fusion,
         GPT-2 Medium MLP at $N{=}2048$, batch $4$.}
\label{tab:epilogue_bw}
\begin{tabular}{lcc}
\toprule
\textbf{Operation} & \textbf{Unfused (GB/step)} & \textbf{Fused (GB/step)} \\
\midrule
FC1 (GEMM)         & 0.84 & 0.84 \\
Bias add            & 0.42 & 0.00 \\
GELU                & 0.42 & 0.00 \\
FC2 (GEMM)         & 0.84 & 0.84 \\
Residual add        & 0.42 & 0.00 \\
\midrule
Total               & 2.94 & 1.68 \\
\bottomrule
\end{tabular}
\end{table}

\section{Methodology}\label{sec:cutlass_method}

\subsection{Kernels Studied}\label{sec:cutlass_method_kernels}
Four shape classes are benchmarked, each one chosen to stress a
different corner of the transformer design space:
\begin{enumerate}
    \item a large square GEMM, $4096 \times 4096 \times 4096$,
          representative of MLP projections at a moderate batch;
    \item a still larger square GEMM, $8192 \times 8192 \times
          8192$, which saturates the Tensor Cores and serves as
          the peak-FLOPS reference;
    \item a skinny GEMM of shape $2048 \times 2048 \times 128$,
          representative of the per-head attention-score and
          attention-output matmuls;
    \item a grouped GEMM for a Mixture-of-Experts layer with $16$
          experts, each of local shape
          $(M_e, N_e, K) = (512, 4096, 4096)$.
\end{enumerate}

\subsection{Implementations Compared}\label{sec:cutlass_method_impl}
Each of these kernels is exercised in three different
implementations: cuBLAS, which is the vendor baseline; CUTLASS~3.x
with tile shapes tuned to the specific problem; and, where
relevant, a hand-rolled CuTe kernel that carries an
epilogue-fused bias+GELU. Across all three the same input tensors
are used---FP16 weights and FP32 accumulators---so that the only
variable is the kernel itself.

\subsection{Metrics}\label{sec:cutlass_method_metrics}
For each kernel we record the sustained throughput in TFLOPS, the
ratio of that throughput against cuBLAS, and the Tensor-Core
utilisation as reported by Nsight Compute. Latencies are taken
from CUDA events after a fixed warm-up of $25$ iterations, and
each configuration is repeated until the relative standard
deviation drops below $1\%$. TFLOPS numbers follow the $2MNK$
convention, with each multiply-add counted as two FLOPs.

\subsection{Hardware and Software}\label{sec:cutlass_method_hw}
All measurements are taken on a single A100-40GB under CUDA~11.8
and CUTLASS~3.2, with the PCIe topology pinned to a single
configuration so that results are reproducible run over run. ECC
is enabled, and the GPU clock is locked to the Tensor-Core boost
frequency to remove clock-frequency variance as a source of noise.

\section{Results}\label{sec:cutlass_numbers}
Table~\ref{tab:cutlass_gemm} summarises the measured throughput
per kernel and implementation.

\begin{table}[h]
\centering
\caption{Representative GEMM performance for LLM shapes.}
\label{tab:cutlass_gemm}
\begin{tabular}{lcccc}
\toprule
\textbf{Shape ($M{\times}N{\times}K$)} & \textbf{Dtype} & \textbf{cuBLAS (TFLOPS)} & \textbf{CUTLASS (TFLOPS)} & \textbf{Ratio} \\
\midrule
$4096 \times 4096 \times 4096$ & FP16 & 280 & 301 & 1.08$\times$ \\
$8192 \times 8192 \times 8192$ & FP16 & 305 & 318 & 1.04$\times$ \\
$2048 \times 2048 \times 128$  & FP16 &  95 & 128 & 1.35$\times$ \\
Grouped (MoE, 16 experts)      & FP16 & 184 & 241 & 1.31$\times$ \\
\bottomrule
\end{tabular}
\end{table}

On the large, cuBLAS-optimal shapes CUTLASS essentially matches
cuBLAS, and on the skinny and grouped shapes---precisely the
patterns that appear in LLM attention and MoE layers---it moves
ahead. The $35\%$ speedup measured on the skinny attention-score
GEMM is traceable directly to a better choice of warp tile,
$64 \times 32 \times 32$, that lines up with the number of
registers available to each warp.

\subsection{Tensor-Core Utilisation}\label{sec:cutlass_tc_util}
Tensor-Core pipeline utilisation, as measured by Nsight Compute's
\texttt{smsp\_\_inst\_executed\_pipe\_tensor} metric, is reported
in Table~\ref{tab:tc_util}. Values above $90\%$ on the large GEMMs
confirm that the kernel is essentially compute-bound; the skinny
GEMM settles at $72\%$ because it has fewer $K$-iterations, and
some of the pipeline stages therefore run empty.

\begin{table}[h]
\centering
\caption{Tensor-Core pipeline utilisation for the CUTLASS
         configurations of Table~\ref{tab:cutlass_gemm}.}
\label{tab:tc_util}
\begin{tabular}{lc}
\toprule
\textbf{Shape} & \textbf{Tensor-Core util. (\%)} \\
\midrule
$4096\times4096\times4096$ & 92.4 \\
$8192\times8192\times8192$ & 95.1 \\
$2048\times2048\times128$  & 72.0 \\
Grouped MoE (16 experts)   & 81.3 \\
\bottomrule
\end{tabular}
\end{table}

\subsection{Epilogue Fusion End-to-End}\label{sec:cutlass_fusion_e2e}
Once the bias+GELU epilogue has been fused into FC1 and the
residual-add epilogue into FC2, the GPT-2 Medium MLP block runs
end-to-end $18\%$ faster at $N = 2048$ with batch~$4$, and $22\%$
faster at $N = 4096$ with batch~$4$. The savings here are
accounted for entirely by the HBM passes eliminated in
Table~\ref{tab:epilogue_bw}.

\section{Summary}\label{sec:cutlass_summary}
The two components introduced in this chapter, CUTLASS and CuTe,
supply the low-level GEMM primitives on which the remainder of
the thesis builds. The Sliding Window Flash Attention kernel of
Chapter~\ref{ch:swfa} and the SGLang serving kernels of
Chapter~\ref{ch:sglang} are both natural clients of these
abstractions. The principal lesson of the chapter is twofold: the
skinny and grouped GEMMs that make up transformer attention and
MoE benefit considerably from a template library that exposes the
tile hierarchy directly to the user, and epilogue fusion recovers
a further $10$--$25\%$ of end-to-end throughput that would
otherwise disappear into HBM round-trips.
